% ============================================================
%  APPENDIX TO CHAPTER 3.  Material NOT in Rudin, imported from
%  Tao I/II, ECNU, USTC, Yu Pin.  Numbered S3.n throughout.
% ============================================================
\rappendix{3}

\begin{roadmap}[How to use this appendix]
Nothing in this appendix is Rudin's text. It is imported from the four
comparison texts and set apart at every level: the running head names the
appendix, section heads carry a reverse-video tag, item numbers are prefixed
\textbf{S}, and every statement names its source.

\medskip
\raggedright
\textbf{Why import anything at all.} Rudin's Chapter 3 is the most complete
chapter in the book, and most of what the other texts add is not missing
theory but missing \emph{leverage}. He defines $\limsup$ as the supremum of
the set of subsequential limits and never gives the formula one actually
computes with. He proves the root and ratio tests and then remarks that both
fail when the limit is $1$ --- without saying what to do next. He proves
summation by parts but extracts no reusable bound from it. Each of those is
one theorem away from being a tool.

\medskip
\textbf{What is here.}

\smallskip
\textbf{A. Upper and lower limits, finished.} The $\inf$-of-$\sup$s formula,
the $\varepsilon$-characterisation, and the algebra --- subadditivity,
products, reciprocals --- that Rudin's 3.19 only begins.

\medskip
\textbf{B. Past the root and ratio tests.} Raabe's test, which compares
against the $p$-series instead of the geometric series; the classical series
that forces it; and the reason the hierarchy of tests can never terminate.

\medskip
\textbf{C. Summation by parts, quantitatively.} Abel's lemma with an explicit
constant, Abel's test with bounded rather than null weights, and the
Dirichlet-kernel estimate that carries Chapter 3 into Fourier series.

\medskip
\textbf{D. Grouping and multiplying.} What inserting parentheses is allowed
to do, and the form of the product theorem that holds for \emph{every}
arrangement of the double array rather than the diagonal one.

\medskip
\textbf{E. Summation without an order.} Rudin's 3.55 says an absolutely
convergent series may be reordered freely. Tao takes that as licence to drop
the order entirely and sum over an index set --- after which double series
and iterated sums are theorems rather than notation.

\medskip
\textbf{How to read it.} Section A pays for itself immediately --- read it
with 3.16--3.19. Sections B--D are best read after 3.44 and 3.50
respectively, and E after 3.55. The exercises are the last section.
\end{roadmap}

% ============================================================
\ssection{Supplement A --- Upper and lower limits, finished}

\begin{roadmap}[Why this supplement belongs to Chapter 3]
\raggedright
\textbf{What Rudin does.} Definition 3.16 puts $E$ for the set of
subsequential limits (in the extended reals) and defines
$\limsup x_n = \sup E$. Theorem 3.17 then proves that $\limsup$ belongs to
$E$ and that no larger number does, and 3.19 gives the one comparison
$a_n \le b_n \Rightarrow \limsup a_n \le \limsup b_n$.

\smallskip
\textbf{The gap.} That definition is superb for proving things and useless
for computing. To find $\limsup$ from it you would first have to find every
subsequential limit --- which is harder than the original problem. The
formula everyone actually uses, $\limsup x_n = \lim_n \sup_{k \ge n} x_k$,
is never stated. Nor is the algebra: no subadditivity, no product rule, no
reciprocal rule.

\smallskip
\textbf{Where the gap bites.} At 3.33 and 3.34, the root and ratio tests,
which are stated in terms of $\limsup$ and therefore cannot be applied until
one can evaluate a $\limsup$; at 3.39, the radius of convergence, where
$1/\limsup\sqrt[n]{|c_n|}$ has to be computed for actual power series; and at
3.37, whose proof compares two upper limits and needs exactly the algebra
that 3.19 does not supply.

\smallskip
\textbf{What this section supplies.} The monotone formula, with the proof
that it agrees with Rudin's definition; the two-clause
$\varepsilon$-characterisation that every later $\limsup$ argument runs on;
and the algebra, including the counterexample showing why the inequalities
are one-directional.
\end{roadmap}

\sitem{S3.1}{Theorem}
\emph{(The formula.) For any real sequence $(x_n)$,}
\[ \limsup_{n\to\infty} x_n = \lim_{n\to\infty}\ \sup_{k \ge n} x_k,
   \qquad
   \liminf_{n\to\infty} x_n = \lim_{n\to\infty}\ \inf_{k \ge n} x_k, \]
\emph{the two limits existing in the extended reals because the sequences
$M_n = \sup_{k\ge n} x_k$ and $m_n = \inf_{k\ge n} x_k$ are monotone.}
\provenance{ECNU \cjk{定理}{Thm} 7.9 (bk.\ p.\,158), which notes that many
texts take this as the \emph{definition}; Rudin never states it.}
\begin{proof}
Since $\{x_k : k \ge n+1\} \subset \{x_k : k \ge n\}$, the suprema satisfy
$M_{n+1} \le M_n$: the sequence $(M_n)$ decreases, so it has a limit $M$ in
the extended reals. Write $E$ for Rudin's set of subsequential limits; we
show $M = \sup E$.

If $L \in E$, say $x_{n_j} \to L$, then $x_{n_j} \le M_n$ for every $j$ with
$n_j \ge n$, so $L \le M_n$ for every $n$ and hence $L \le M$. Thus
$\sup E \le M$.

Conversely $M \in E$. Choose indices $n_1 < n_2 < \cdots$ so that
$x_{n_k} > M_{n_k} - 1/k$, which is possible because $M_{n_k}$ is the
supremum of the tail beginning at $n_k$. Since also $x_{n_k} \le M_{n_k}$
and $M_{n_k} \to M$, we get $x_{n_k} \to M$. Hence $M \le \sup E$.
\end{proof}

\begin{proofwalk}[How the proof flows]
  \item \textbf{Choose the method.} Two descriptions of one number are
        given, so the proof is two inequalities. Neither side can be
        computed, so each must be pinned by an approximation argument.
  \item \textbf{Get monotonicity free.} Discarding terms can only lower a
        supremum. That single observation is what makes $\lim M_n$ exist at
        all --- by monotone convergence (3.14), with no boundedness needed
        once $\pm\infty$ are allowed.
  \item \textbf{One inequality by domination.} Every subsequential limit is
        trapped below every tail supremum, so $\sup E \le M$. Nothing is
        constructed here.
  \item \textbf{The other by construction.} To show $M \le \sup E$ it is not
        enough that $M$ bounds things --- $M$ must \emph{be} a subsequential
        limit, so a subsequence converging to it has to be built. The
        recipe is the defining property of a supremum: at stage $k$ pick a
        term within $1/k$ of $M_{n_k}$.
  \item \textbf{Audit.} The least-upper-bound property is spent once, in
        step 4, to produce each $x_{n_k}$; monotone convergence once, in
        step 2. Note which clause of Definition 1.8 does which job: ``upper
        bound'' gives step 3, ``least'' gives step 4 --- the same division of
        labour as in Rudin's 1.21.
\end{proofwalk}

\begin{unpack}[Why this is the version you compute with]
$M_n$ is a decreasing sequence --- taking the supremum over a smaller tail
cannot increase it --- and it is bounded below whenever $(x_n)$ is. So the
limit exists by monotone convergence, with no reference to subsequences at
all. To evaluate $\limsup$ you now do two things you know how to do: take a
supremum, then take a monotone limit.

Rudin's $E$-based definition and this one agree, and the agreement is worth
seeing as two inequalities. Every subsequential limit is $\le M_n$ for all
$n$, hence $\le \lim M_n$; and $\lim M_n$ is itself a subsequential limit,
because for each $n$ one can pick an index $k_n \ge n$ with
$x_{k_n} > M_n - 1/n$. The second half is where the work is, and it is the
same construction as ECNU's proof that the largest cluster point exists.
\end{unpack}

\begin{center}
\begin{tikzpicture}[x=1mm, y=1mm]
  \draw[axis] (0,0) -- (100,0);
  \draw[axis] (0,0) -- (0,42);
  \draw[guide] (2,26) -- (95,26);
  \draw[guide] (2,10) -- (95,10);
  \foreach \n/\v in {1/34, 2/6, 3/30, 4/8, 5/28, 6/9, 7/27, 8/9.5,
                     9/26.5, 10/9.8, 11/26.2, 12/9.9, 13/26.1, 14/9.95}{
    \node[ptfill, minimum size=2.4pt] at ({6.4*\n},\v) {};}
  % sup staircase, with risers
  \draw[thin] (6.4,34) -- (19.2,34) -- (19.2,30) -- (32,30) -- (32,28)
    -- (44.8,28) -- (44.8,27) -- (57.6,27) -- (57.6,26.5) -- (70.4,26.5)
    -- (70.4,26.2) -- (89.6,26.2);
  % inf staircase, with risers
  \draw[thin] (12.8,6) -- (25.6,6) -- (25.6,8) -- (38.4,8) -- (38.4,9)
    -- (51.2,9) -- (51.2,9.5) -- (64,9.5) -- (64,9.8) -- (76.8,9.8)
    -- (76.8,9.9) -- (89.6,9.9);
  \node[mlbl, anchor=west] at (96,26) {$\limsup$};
  \node[mlbl, anchor=west] at (96,10) {$\liminf$};
  \node[lbl, anchor=west] at (30,38) {$M_n=\sup_{k\ge n}x_k$ falls};
  \node[lbl, anchor=west] at (30,2.6) {$m_n=\inf_{k\ge n}x_k$ rises};
\end{tikzpicture}
\end{center}
\figcap{The two staircases. Discarding a term can only lower a supremum and
raise an infimum, so $M_n$ decreases and $m_n$ increases; their limits are
the upper and lower limits. The picture also shows why the sequence converges
exactly when the staircases meet.}

\sitem{S3.2}{Theorem}
\emph{($\varepsilon$-characterisation.) $\bar A = \limsup x_n$ if and only if
for every $\varepsilon>0$: (i) $x_n < \bar A + \varepsilon$ for all
sufficiently large $n$, and (ii) $x_n > \bar A - \varepsilon$ for infinitely
many $n$. Equivalently: every $\alpha > \bar A$ is exceeded only finitely
often, and every $\beta < \bar A$ is exceeded infinitely often.}
\provenance{ECNU \cjk{定理}{Thm} 7.7 and 7.7$'$ (bk.\ p.\,157). Rudin's
3.17(b) is condition (i) alone.}
\begin{proof}
Let $\bar A = \limsup x_n$ and use S3.1. Given $\varepsilon>0$, the numbers
$M_n$ decrease to $\bar A$, so $M_N < \bar A + \varepsilon$ for some $N$, and
then $x_n \le M_N < \bar A+\varepsilon$ for all $n \ge N$: that is (i). And
$M_n \ge \bar A$ for every $n$, so $\sup_{k \ge n} x_k > \bar A - \varepsilon$,
which produces some $k \ge n$ with $x_k > \bar A-\varepsilon$; as $n$ is
arbitrary there are infinitely many such $k$, which is (ii).

Conversely, suppose $A$ satisfies (i) and (ii). By (i), $M_N \le
A+\varepsilon$ for some $N$, so $\bar A \le A + \varepsilon$; by (ii),
$M_n > A - \varepsilon$ for every $n$, so $\bar A \ge A - \varepsilon$. As
$\varepsilon$ was arbitrary, $\bar A = A$ by the $\varepsilon$-principle
(S1.9).
\end{proof}

\noindent
Rudin's 3.17 says $\limsup$ belongs to $E$ and that nothing above it does.
The $\varepsilon$-form is the same content converted into the two statements
one actually verifies.\kn{Read the pair as ``eventually below, infinitely
often nearly above.'' Almost every $\limsup$ argument in later chapters ---
the root test at 3.33, the radius of convergence at 3.39 --- uses exactly
these two clauses, one to get an upper bound that eventually holds and one to
find infinitely many terms witnessing sharpness.}

\sitem{S3.3}{Theorem}
\emph{(Algebra of upper and lower limits.) For bounded real sequences,}
\begin{enumerate}[label=(\alph*), leftmargin=2.2em, itemsep=2pt, topsep=3pt]
  \item $\liminf x_n = -\limsup(-x_n)$;
  \item $\limsup(x_n+y_n) \le \limsup x_n + \limsup y_n$, \emph{with
        equality whenever $\lim x_n$ exists;}
  \item $\liminf x_n + \liminf y_n \le \liminf(x_n+y_n)$;
  \item \emph{if $x_n, y_n > 0$, then}
        $(\liminf x_n)(\liminf y_n) \le \liminf(x_ny_n)$;
  \item \emph{if $\liminf x_n > 0$, then}
        $\limsup(1/x_n) = 1/\liminf x_n$.
\end{enumerate}
\provenance{ECNU \cjk{定理}{Thm} 7.8, \cjk{例}{Ex} 2, and
\cjk{习题}{Ex} 7.2 \#2 (bk.\ pp.\,158--159).}
\begin{proof}
All five follow from S3.1 together with elementary facts about suprema over a
common tail.

(a) $\inf_{k\ge n}(-x_k) = -\sup_{k \ge n} x_k$; let $n \to \infty$.

(b) $\sup_{k\ge n}(x_k+y_k) \le \sup_{k\ge n}x_k + \sup_{k\ge n}y_k$, since
each term on the left is dominated by the right side; let $n \to \infty$. If
$x_n \to x$ then applying the inequality to $y_n = (x_n+y_n) + (-x_n)$ gives
$\limsup y_n \le \limsup(x_n+y_n) + \limsup(-x_n) = \limsup(x_n+y_n) - x$,
which is the reverse inequality.

(c) Apply (b) to $-x_n$ and $-y_n$ and use (a).

(d) For positive terms, $\inf_{k\ge n}(x_ky_k) \ge
(\inf_{k\ge n}x_k)(\inf_{k\ge n}y_k)$; let $n\to\infty$.

(e) If $\liminf x_n > 0$ then $x_k > 0$ for large $k$ and
$\sup_{k\ge n}(1/x_k) = 1/\inf_{k\ge n}x_k$; let $n \to \infty$.
\end{proof}

\begin{pitfall}[The inequalities are one-directional, and that is the point]
It is tempting to expect $\limsup(x_n+y_n) = \limsup x_n + \limsup y_n$. Take
\[ x_n = (-1)^n, \qquad y_n = (-1)^{n+1} . \]
Then $x_n + y_n = 0$ for every $n$, so the left side is $0$, while the right
side is $1+1 = 2$. The two sequences achieve their upper limits on
\emph{disjoint} sets of indices, and $\limsup$ cannot see that they never
peak together.

Equality is restored the moment one of the sequences converges, because a
convergent sequence peaks everywhere at once. That is why Rudin can get away
with 3.19 alone: he never adds two genuinely oscillating sequences. When you
do --- in the root test applied to a product, say --- the inequality is all
you have.
\end{pitfall}

% ============================================================
\ssection{Supplement B --- Past the root and ratio tests}

\begin{roadmap}[Why this supplement belongs to Chapter 3]
\raggedright
\textbf{What Rudin does.} Theorems 3.33 and 3.34 give the root and ratio
tests; the Remark after 3.35 observes that both are silent when the relevant
limit equals $1$, illustrates with $\Sigma 1/n$ and $\Sigma 1/n^2$, and
moves on.

\smallskip
\textbf{The gap.} The remark diagnoses a failure without saying what to do
about it, and without explaining \emph{why} the failure is unavoidable. A
reader is left with the impression that the borderline is a small blind spot
rather than the entrance to an infinite hierarchy.

\smallskip
\textbf{Where the gap bites.} On every series that matters at the boundary.
Rudin's own 3.28 --- $\Sigma 1/n^p$ --- sits exactly there, and he must
reach it by a separate device, the condensation test 3.27. The
hypergeometric series, the central example of nineteenth-century analysis,
is invisible to both of his tests. And at 3.39 the radius of convergence of
a power series is often decided precisely on the circle, where root and ratio
give nothing.

\smallskip
\textbf{What this section supplies.} Raabe's test, which replaces the
geometric yardstick by the $p$-series one and so resolves the whole borderline
band; the classical series that forces it; and the reason no test can be
last --- there is no slowest convergent series, so Rudin's remark is a
structural fact rather than an omission.
\end{roadmap}

\sitem{S3.4}{Theorem}
\emph{(Raabe's test.) Let $u_n > 0$. If there is $r > 1$ with}
\[ n\Bigl(1 - \frac{u_{n+1}}{u_n}\Bigr) \ge r \qquad (n > N_0), \]
\emph{then $\Sigma u_n$ converges. If $n\bigl(1 - u_{n+1}/u_n\bigr) \le 1$
for all $n > N_0$, then $\Sigma u_n$ diverges. In particular, if the limit
$r = \lim n\bigl(1-u_{n+1}/u_n\bigr)$ exists, then $r>1$ gives convergence
and $r<1$ divergence.}
\provenance{ECNU \cjk{拉贝判别法}{Raabe} \cjk{定理}{Thm} 12.10
(\cjk{下}{vol.2} bk.\ pp.\,14--15).}
\begin{proof}
\emph{Convergence.} Suppose $n(1-u_{n+1}/u_n) \ge r > 1$ for $n>N_0$, so
that $u_{n+1}/u_n \le 1 - r/n$. Fix $p$ with $1<p<r$. Bernoulli's inequality
$(1+t)^p \ge 1+pt$, applied with $t=-1/n$, gives
\[ \Bigl(\frac{n-1}{n}\Bigr)^{p} \;\ge\; 1 - \frac{p}{n} \;\ge\;
   1 - \frac{r}{n} \;\ge\; \frac{u_{n+1}}{u_n} , \]
that is $n^p u_{n+1} \le (n-1)^p u_n$. So the numbers $(n-1)^pu_n$ decrease
for $n>N_0$ and are therefore bounded by some $C$, giving
$u_n \le C/(n-1)^p$. Since $p>1$ the series $\Sigma (n-1)^{-p}$ converges
(3.28), and the comparison test 3.25 finishes.

\emph{Divergence.} If $n(1-u_{n+1}/u_n) \le 1$ then
$u_{n+1}/u_n \ge 1-1/n = (n-1)/n$, so $nu_{n+1} \ge (n-1)u_n$ and the numbers
$(n-1)u_n$ increase; hence $(n-1)u_n \ge c>0$ and $u_n \ge c/(n-1)$. Since
$\Sigma 1/(n-1)$ diverges, so does $\Sigma u_n$.
\end{proof}

\begin{proofwalk}[How the proof flows]
  \item \textbf{Choose the method.} The ratio test failed because it compares
        with a geometric series. So compare with a $p$-series instead --- but
        the hypothesis is about \emph{ratios}, so the comparison must be made
        between ratios rather than between terms.
  \item \textbf{Compute the target ratio.} For $v_n = n^{-p}$ the ratio is
        $v_{n+1}/v_n = \bigl((n)/(n+1)\bigr)^{p}$, which Bernoulli bounds
        below by $1-p/n$. So ``$u_{n+1}/u_n \le 1 - p/n$'' is exactly ``the
        terms decay at least as fast as $n^{-p}$.''
  \item \textbf{Buy room with $p<r$.} The hypothesis gives $r$; the
        comparison needs a strict $p$ above $1$. Any $p$ between them works,
        and this is the only place the strictness $r>1$ is used.
  \item \textbf{Convert to a monotone quantity.} The ratio inequality says
        precisely that $(n-1)^pu_n$ decreases. A decreasing positive sequence
        is bounded, and boundedness is what the comparison test wants.
  \item \textbf{Audit.} Bernoulli is spent once, in step 2; the strict
        inequality $r>1$ once, in step 3; the convergence of the $p$-series
        once, at the end. Run the same four steps with the inequalities
        reversed and $p=1$ and the divergence half falls out --- which is why
        the borderline case $r=1$ is settled by neither.
\end{proofwalk}

\begin{unpack}[What changed: the yardstick]
The root and ratio tests compare $\Sigma u_n$ against a \emph{geometric}
series. When $u_{n+1}/u_n \to 1$ the comparison is vacuous, because the
geometric series with ratio $1$ is exactly the borderline.

Raabe compares against the \emph{$p$-series} instead. Writing
$u_{n+1}/u_n \approx 1 - r/n$ and comparing with $(1-1/n)^r$ shows that the
quantity $n(1 - u_{n+1}/u_n)$ measures which power of $n$ the terms decay
like: $r > 1$ means faster than $1/n$, and $r < 1$ slower. So the test is
``compare with $\Sigma n^{-r}$'' with the comparison done through ratios
rather than the terms themselves.

That reframing tells you immediately what the next rung must be. Since the
$p$-series is itself borderline at $p=1$, Raabe fails when $r=1$ --- and the
next test compares against $\Sigma 1/\bigl(n(\log n)^{p}\bigr)$, which
Rudin's own 3.29 has already introduced.
\end{unpack}

\begin{deeper}[The series that forces the issue, and why the ladder never ends]
The standard example is
\[ \sum_{n=1}^{\infty}
   \Bigl[\frac{1 \cdot 3 \cdot 5 \cdots (2n-1)}
              {2 \cdot 4 \cdot 6 \cdots (2n)}\Bigr]^{s}, \]
whose ratio tends to $1$ for \emph{every} $s$, so ratio and root say nothing
whatever. Raabe gives $r = s/2$, hence divergence for $s = 1$ and $s = 2$ and
convergence for $s = 3$. The same computation settles the hypergeometric
series $F(\alpha,\beta,\gamma,1)$, which converges precisely when
$\gamma > \alpha + \beta$ --- historically the problem the whole hierarchy of
tests was invented for.

ECNU then states the moral, and it is worth taking seriously: one can always
build a finer test, and the process never terminates, because \emph{there is
no slowest convergent series}. Given any convergent $\Sigma a_n$ one can
manufacture a convergent series whose terms dominate it eventually, so no
single comparison series can serve as a universal boundary. Rudin's 3.35
remark --- that when the ratio is $1$ ``the series may converge or diverge''
--- is therefore not laziness but a structural fact.
\end{deeper}

\begin{deeper}[The integral test, and why Rudin cannot state it here]
Both Chinese texts prove: if $f$ is nonnegative and decreasing on
$[1,\infty)$, then $\Sigma f(n)$ and $\int_1^{\infty} f(x)\,dx$ converge
together, via the sandwich
$\sum_{k=2}^{n+1} f(k) \le \int_1^{n+1} f \le \sum_{k=1}^{n} f(k)$.

It is the fastest route to the $p$-series and to the iterated-logarithm
ladder $\Sigma 1/\bigl(n\log n (\log\log n)^p\bigr)$. Rudin gets both by
\emph{Cauchy condensation} (3.27--3.29) instead, and he has no choice: the
integral does not exist until Chapter 6, so a Chapter 3 integral test would
be circular in his development. Come back to this box after 6.13, where the
comparison becomes available --- and notice that condensation and the
integral test are the same idea, one summing over dyadic blocks and the other
over unit intervals.
\end{deeper}

% ============================================================
\ssection{Supplement C --- Summation by parts, quantitatively}

\begin{roadmap}[Why this supplement belongs to Chapter 3]
\raggedright
\textbf{What Rudin does.} Theorem 3.41 is the partial summation formula, and
he uses it once, immediately, to prove 3.42 --- if the partial sums of
$\Sigma a_n$ are bounded and $b_n \downarrow 0$ then $\Sigma a_nb_n$
converges --- from which 3.43 (alternating series) and 3.44 follow.

\smallskip
\textbf{The gap.} The formula is consumed on the spot. No \emph{bound} is
extracted from it, and no second theorem: Rudin has the Dirichlet test but
not the Abel test, which trades the hypotheses the other way round.

\smallskip
\textbf{Where the gap bites.} At 3.44, where $\Sigma z^n/n^{\alpha}$ on
$|z|=1$ is handled by the Dirichlet route because the Abel route is
unavailable; and, more seriously, in Chapter 7, where the uniform versions of
these tests are wanted and a uniform conclusion needs an explicit constant,
not merely a convergence statement.

\smallskip
\textbf{What this section supplies.} The bound --- Abel's lemma, with the
constant $3$ --- from which Rudin's 3.42 and its missing twin both follow in
a line; and the Dirichlet kernel estimate, which is the standard bridge from
this chapter to Fourier series and which 3.44 stops just short of.
\end{roadmap}

\sitem{S3.5}{Theorem}
\emph{(Abel's lemma.) Let $(\varepsilon_k)$ be monotone, let
$\sigma_k = v_1+\cdots+v_k$, and suppose $|\sigma_k| \le A$ for
$k \le n$. Put $\varepsilon = \max_k |\varepsilon_k|$. Then}
\[ \Bigl|\sum_{k=1}^{n} \varepsilon_k v_k\Bigr| \le 3\varepsilon A . \]
\provenance{ECNU \cjk{阿贝尔引理}{Abel's lemma} (\cjk{下}{vol.2}
bk.\ pp.\,21--22).}
\begin{proof}
Partial summation (3.41) gives
\[ \sum_{k=1}^{n} \varepsilon_k v_k
   = \sum_{k=1}^{n-1} (\varepsilon_k - \varepsilon_{k+1})\sigma_k
     + \varepsilon_n \sigma_n . \]
Because $(\varepsilon_k)$ is monotone, the differences
$\varepsilon_k - \varepsilon_{k+1}$ all have one sign, so
\[ \sum_{k=1}^{n-1} |\varepsilon_k-\varepsilon_{k+1}|
   = |\varepsilon_1 - \varepsilon_n| \le 2\varepsilon . \]
Bounding every $|\sigma_k|$ by $A$ therefore gives
$\bigl|\sum \varepsilon_kv_k\bigr| \le 2\varepsilon A + \varepsilon A
= 3\varepsilon A$.
\end{proof}

\noindent
The proof is Rudin's 3.41 followed by one observation: because
$(\varepsilon_k)$ is monotone, the differences
$\varepsilon_k - \varepsilon_{k+1}$ all have the same sign, so
$\sum_k |\varepsilon_k - \varepsilon_{k+1}| = |\varepsilon_1 -
\varepsilon_n| \le 2\varepsilon$, and the boundary terms contribute
$\varepsilon A$ more.\kn{This is exactly the telescope Rudin performs inside
the proof of 3.42, extracted and given a name. Rudin's version is consumed
on the spot to prove one convergence theorem; as a standalone inequality it
is reusable, and it is what the uniform-convergence versions in Chapter 7
run on.}

\sitem{S3.6}{Theorem}
\emph{(Abel's test.) If $(a_n)$ is monotone and \textbf{bounded} --- not
necessarily null --- and $\Sigma b_n$ converges, then $\Sigma a_nb_n$
converges.}
\provenance{ECNU \cjk{定理}{Thm} 12.15 (\cjk{下}{vol.2} bk.\ p.\,22).}
\begin{proof}
Being monotone and bounded, $(a_n)$ converges, say to $a$; put
$b_n = a_n - a$, so $(b_n)$ is monotone with $b_n \to 0$. Then
\[ \sum a_nc_n = a\sum c_n + \sum b_nc_n . \]
The first series converges by hypothesis. For the second, the partial sums of
$\Sigma c_n$ are bounded --- a convergent sequence is bounded --- and
$b_n \downarrow 0$, so Rudin's 3.42 applies. Hence $\Sigma a_nc_n$
converges.
\end{proof}

\begin{deeper}[Two tests, and which hypothesis each spends]
Rudin's 3.42 is the \emph{Dirichlet} test: partial sums of $\Sigma b_n$
merely bounded, but $a_n \downarrow 0$. Abel's test trades those hypotheses:
$\Sigma b_n$ must actually converge, but $(a_n)$ need only be monotone and
bounded. Neither implies the other, and both fall out of Abel's lemma.

\smallskip
\begin{tabular}{@{}p{0.3\linewidth}p{0.3\linewidth}p{0.3\linewidth}@{}}
\toprule
& \textbf{Dirichlet (Rudin 3.42)} & \textbf{Abel (S3.6)} \\
\midrule
on $\Sigma b_n$   & partial sums bounded & converges \\
on $(a_n)$        & monotone, $\to 0$    & monotone, bounded \\
typical use       & $\Sigma \frac{\sin nx}{n}$ & multiplying a convergent
series by a damping factor \\
\bottomrule
\end{tabular}

\smallskip
Rudin never states Abel's form, and its absence is felt at 3.44, where he
handles $\Sigma z^n/n^{\alpha}$ on $|z|=1$ by the Dirichlet route.
\end{deeper}

\sitem{S3.7}{Theorem}
\emph{(The Dirichlet kernel bound.) For $x \in (0,2\pi)$,}
\[ \frac12 + \sum_{k=1}^{n}\cos kx
   = \frac{\sin\bigl((n+\tfrac12)x\bigr)}{2\sin(x/2)} , \]
\emph{so the partial sums of $\Sigma\cos kx$ are bounded on any interval
$[\delta, 2\pi-\delta]$. Consequently, if $a_1 \ge a_2 \ge \cdots$ with
$a_n \to 0$, then $\Sigma a_n \cos nx$ and $\Sigma a_n \sin nx$ converge for
every $x \in (0,2\pi)$.}
\provenance{ECNU \cjk{例}{Ex} 3 (\cjk{下}{vol.2} bk.\ pp.\,22--23).}
\begin{proof}
For $x \in (0,2\pi)$ we have $\sin(x/2) \ne 0$. The product formula
$2\cos kx\,\sin(x/2) = \sin\bigl((k+\tfrac12)x\bigr) -
\sin\bigl((k-\tfrac12)x\bigr)$ makes the sum telescope:
\[ 2\sin\tfrac{x}{2}\Bigl[\tfrac12 + \sum_{k=1}^{n}\cos kx\Bigr]
   = \sin\tfrac{x}{2} + \sin\bigl((n+\tfrac12)x\bigr) - \sin\tfrac{x}{2}
   = \sin\bigl((n+\tfrac12)x\bigr), \]
which is the identity. Hence
$\bigl|\sum_{k\le n}\cos kx\bigr| \le 1/\bigl(2|\sin(x/2)|\bigr) + \tfrac12$,
a bound independent of $n$ and finite for each fixed
$x \in (0,2\pi)$. With $a_n \downarrow 0$, Rudin's 3.42 gives convergence of
$\Sigma a_n\cos nx$; the same telescoping with $\sin$ in place of $\cos$
bounds the partial sums of $\Sigma \sin kx$ and settles that series too.
\end{proof}

\noindent
This is the standard bridge from Chapter 3 to Fourier series: it makes
$\Sigma (\sin nx)/n$ and $\Sigma (\cos nx)/n$ convergent everywhere on
$(0,2\pi)$, none of which follows from absolute convergence, since
$\Sigma 1/n$ diverges.\kn{The identity is proved by multiplying the left side
by $2\sin(x/2)$ and telescoping the products
$2\cos kx\sin(x/2) = \sin((k+\frac12)x) - \sin((k-\frac12)x)$. Note where the
bound degrades: the denominator $2\sin(x/2)$ tends to $0$ as $x \to 0$, so
the partial sums are bounded on each $[\delta,2\pi-\delta]$ but not
uniformly on $(0,2\pi)$ --- which is exactly the phenomenon Chapter 7 will
call non-uniform convergence.}

% ============================================================
\ssection{Supplement D --- Grouping and multiplying}

\begin{roadmap}[Why this supplement belongs to Chapter 3]
\raggedright
\textbf{What Rudin does.} Definition 3.48 defines the product of two series
by the anti-diagonals $c_n = \sum_k a_kb_{n-k}$, motivated by power series;
3.49 shows the product of two conditionally convergent series can diverge;
3.50 (Mertens) proves convergence to $AB$ when one factor is absolutely
convergent.

\smallskip
\textbf{The gap.} Two operations on series are performed silently throughout
the chapter and never licensed. \emph{Grouping} --- inserting parentheses ---
is used inside the condensation proof at 3.27 and in the Remark after 3.44,
but never stated as a theorem, so a reader has no way to know it is legal
while rearrangement is not. And the choice of the anti-diagonal ordering at
3.48 is presented as forced by power series, leaving open what happens under
any other arrangement of the same products.

\smallskip
\textbf{Where the gap bites.} At 3.49 especially, where the divergence of a
product looks like a pathology of multiplication rather than the price of
insisting on one particular traversal of a doubly-indexed array.

\smallskip
\textbf{What this section supplies.} The grouping theorem with its
one-directional caveat, so that the legal manipulation is cleanly separated
from the illegal one; and Cauchy's theorem, that under absolute convergence
\emph{every} arrangement of the array converges to $AB$ --- which puts 3.50
and 3.49 in their proper places.
\end{roadmap}

\sitem{S3.8}{Theorem}
\emph{(Grouping.) Inserting parentheses in a convergent series ---
that is, passing to any subsequence of the partial sums by grouping
consecutive terms --- changes neither convergence nor the sum.}
\provenance{ECNU \cjk{定理}{Thm} 12.4 (\cjk{下}{vol.2} bk.\ pp.\,4--5).}
\begin{proof}
Let $s_n$ be the partial sums of the original series and let the parentheses
end at indices $n_1 < n_2 < \cdots$. The partial sums of the grouped series
are exactly $s_{n_1}, s_{n_2}, \dots$ --- a subsequence of $(s_n)$. A
subsequence of a convergent sequence converges to the same limit (3.6(a)).
\end{proof}

\begin{pitfall}[Grouping is one-directional, unlike rearrangement]
The grouped partial sums form a \emph{subsequence} of $(s_n)$, and a
subsequence of a convergent sequence converges to the same limit --- that is
the whole proof. The converse fails:
\[ (1-1) + (1-1) + (1-1) + \cdots = 0 \]
converges while $1 - 1 + 1 - 1 + \cdots$ diverges. So parentheses may be
\emph{inserted} into a convergent series but not \emph{removed} from one.

Keep this distinct from the rearrangement theorem 3.54. Rearranging changes
the order of the terms and can change the sum of a conditionally convergent
series to anything at all; grouping preserves the order and is always safe in
the direction stated. Rudin never states the grouping theorem, although he
uses it --- in the condensation proof at 3.27 and in the Remark after 3.44.
\end{pitfall}

\sitem{S3.9}{Theorem}
\emph{(Cauchy's product theorem, any arrangement.) If $\Sigma u_n = A$ and
$\Sigma v_n = B$ both converge absolutely, then the series formed from all
products $u_iv_j$, each taken once, \textbf{in any order whatever},
converges absolutely to $AB$.}
\provenance{ECNU \cjk{柯西定理}{Cauchy's thm} 12.14 (\cjk{下}{vol.2}
bk.\ pp.\,19--21).}
\begin{proof}
Let $(w_k)$ enumerate the products $u_iv_j$ in any order. Any finite
collection of them involves indices $i,j \le N$ for some $N$, so
\[ \sum_{k \le K} |w_k| \;\le\;
   \Bigl(\sum_{i\le N}|u_i|\Bigr)\Bigl(\sum_{j\le N}|v_j|\Bigr)
   \;\le\; \Bigl(\sum_i |u_i|\Bigr)\Bigl(\sum_j |v_j|\Bigr) . \]
The partial sums of $\Sigma|w_k|$ are bounded, so $\Sigma w_k$ converges
absolutely; by 3.55 its sum is the same for every enumeration. Evaluate it
along the square ordering, in which the partial sum after the $N$th shell is
exactly $A_NB_N$, and $A_NB_N \to AB$ by 3.3(c).
\end{proof}

\begin{center}
\begin{tikzpicture}[x=1mm, y=1mm]
  % square / shell ordering
  \foreach \i in {0,7,...,28}{\foreach \j in {0,7,...,28}{
    \node[ptfill, minimum size=1.6pt] at (\i,\j) {};}}
  \node[ptopen] at (0,0) {};
  \draw[rmarrow] (7,-1.4) -- (7,7) -- (-1.4,7);
  \draw[rmarrow] (14,-1.4) -- (14,14) -- (-1.4,14);
  \draw[rmarrow] (21,-1.4) -- (21,21) -- (-1.4,21);
  \node[lbl, anchor=north] at (14,-4.5) {square (shell) order};
  \node[mlbl, anchor=north east] at (-1,-1) {$u_1v_1$};
  % diagonal ordering
  \begin{scope}[xshift=58mm]
  \foreach \i in {0,7,...,28}{\foreach \j in {0,7,...,28}{
    \node[ptfill, minimum size=1.6pt] at (\i,\j) {};}}
  \foreach \d in {7,14,21,28}{
    \draw[rmarrow] (0,\d) -- (\d,0);}
  \node[lbl, anchor=north] at (14,-4.5) {diagonal order (Rudin's $c_n$)};
  \node[mlbl, anchor=north east] at (-1,-1) {$u_1v_1$};
  \end{scope}
\end{tikzpicture}
\end{center}
\figcap{Two ways to enumerate the same array of products $u_iv_j$. Along
shells the partial sums are exactly $A_nB_n$, which is why the square order
proves the theorem; along anti-diagonals they are Rudin's $C_n$, which is
the order that makes the product a \emph{series} in the sense of 3.48.
Absolute convergence is what makes the two agree.}

\begin{deeper}[How this relates to Mertens (3.50)]
The two theorems point in opposite directions. Rudin's 3.50 fixes the
diagonal ordering --- the one that matters for power series --- and weakens
the hypothesis: only \emph{one} factor need converge absolutely. S3.9 keeps
the stronger hypothesis, both absolutely convergent, and buys freedom of
arrangement.

The proof of S3.9 is the cleaner of the two, because the square ordering
makes the partial sums \emph{exactly} $A_nB_n$ rather than approximately so;
all the difficulty in Mertens' theorem comes from the diagonal ordering
cutting the array along a line the factors do not respect. Seen this way,
Example 3.49 --- where the Cauchy product of two conditionally convergent
series diverges --- is not a pathology but the expected price of insisting on
the diagonal.
\end{deeper}

% ============================================================
\ssection{Supplement E --- Summation without an order}

\begin{roadmap}[Why this supplement belongs to Chapter 3]
\raggedright
\textbf{What Rudin does.} Theorem 3.54 shows a conditionally convergent
series can be rearranged to any prescribed value, and 3.55 shows an
absolutely convergent one cannot: every rearrangement has the same sum.

\smallskip
\textbf{The gap.} 3.55 is stated as a stability property --- the sum survives
reordering --- rather than as what it actually is: a licence to forget the
ordering altogether. Rudin never draws that conclusion, so a series remains,
for him, a function on $\N$, and anything indexed differently must be
converted to one by hand.

\smallskip
\textbf{Where the gap bites.} At 3.48, where the product of two series has to
be \emph{defined} by choosing a traversal of a square array, and at 3.49,
where the choice turns out to matter. Later, in Chapter 8, power series are
multiplied and rearranged repeatedly, and each manipulation is justified
afresh.

\smallskip
\textbf{What this section supplies.} Tao's reading of 3.55: for absolutely
summable families the sum is a function of the \emph{index set}, so it may be
defined with no order at all. Double series and iterated sums then become
theorems --- Fubini for series --- rather than notational accidents.
\end{roadmap}

\sitem{S3.10}{Definition}
\emph{(Sums over arbitrary index sets.) Let $X$ be a set and
$f : X \to \R$. If $X$ is countable, choose a bijection $g : \N \to X$ and
put $\sum_{x \in X} f(x) = \sum_{n=0}^{\infty} f(g(n))$, provided the latter
converges absolutely; the value does not depend on $g$. If $X$ is
uncountable, the same definition applies once one knows that
$\{x : f(x) \ne 0\}$ is at most countable.}
\provenance{Tao I, \S8.2, Def.\ 8.2.1--8.2.4 and Lemma 8.2.5
(bk.\ pp.\,187--189).}

\begin{unpack}[What Rudin's 3.55 becomes]
Rudin proves that an absolutely convergent series may be rearranged freely
(3.55) and that a conditionally convergent one may be rearranged to any sum
whatever (3.54). Tao reads the first as a licence: since the value is
independent of the enumeration, \emph{the ordering was never part of the
data}. A series is then a function on an index set, and $\sum_{x\in X}f(x)$
is well defined with no order in sight.

The dividend is that families indexed by things other than $\N$ --- pairs,
grids, arbitrary countable sets --- need no separate theory. Rudin has to
manufacture an ordering to define the Cauchy product at 3.48 (his $c_n$ are
the anti-diagonals), and Example 3.49's failure is precisely the price of
that manufactured ordering.
\end{unpack}

\sitem{S3.11}{Lemma}
\emph{If $\sum_{x\in X} f(x)$ is absolutely convergent, then
$\{x \in X : f(x) \ne 0\}$ is at most countable --- however large $X$ is.}
\provenance{Tao I, Lemma 8.2.5 (bk.\ p.\,188).}

\noindent
So an absolutely summable family over an uncountable index set is really a
series in disguise.\kn{The proof is the standard countability argument in
Rudin's own idiom: for each $n$ the set where $|f| > 1/n$ must be finite,
or the sum would exceed any bound; a countable union of finite sets is
countable. Compare Exercise S2.18 in the Chapter 2 appendix --- disjoint
intervals are countable for the same reason.}

\sitem{S3.12}{Theorem}
\emph{(Fubini for series.) If $f : \N \times \N \to \R$ has
$\sum_{(n,m)} f(n,m)$ absolutely convergent, then}
\[ \sum_{n=0}^{\infty}\Bigl(\sum_{m=0}^{\infty} f(n,m)\Bigr)
   = \sum_{(n,m) \in \N\times\N} f(n,m)
   = \sum_{m=0}^{\infty}\Bigl(\sum_{n=0}^{\infty} f(n,m)\Bigr). \]
\provenance{Tao I, Thm 8.2.2 (bk.\ p.\,188). Rudin states nothing about
iterated sums in Chapter 3.}
\begin{proof}
Write $S$ for the unordered sum, which exists by S3.10. Given
$\varepsilon>0$, absolute convergence supplies a finite set
$F \subset \N\times\N$ with $\sum_{(n,m)\notin F} |f(n,m)| < \varepsilon$;
choose $N$ with $F \subset [0,N]^2$.

Each inner sum $g(n) = \sum_m f(n,m)$ converges absolutely, and
$\sum_{n>N} |g(n)| \le \sum_{n>N}\sum_m |f(n,m)| < \varepsilon$. Comparing
$\sum_{n \le N} g(n)$ with $S$, the terms they fail to share all lie outside
$F$, so the two differ by at most $\varepsilon$. Hence
$\bigl|\sum_{n\le N}g(n) - S\bigr| \le \varepsilon$ and, adding the tail
estimate, $\bigl|\sum_{n} g(n) - S\bigr| \le 2\varepsilon$. As $\varepsilon$
was arbitrary, $\sum_n g(n) = S$; the argument with the roles of $n$ and $m$
exchanged gives the other iterated sum.
\end{proof}

\begin{deeper}[Why this is the theorem the Cauchy product wanted]
S3.9 said that an absolutely convergent double array may be summed in any
arrangement. S3.12 says the two most useful arrangements --- by rows and by
columns --- also agree with each other and with the unordered sum. Together
they are the complete story of the array $u_iv_j$, of which Rudin's 3.50 is
the single diagonal slice.

Absolute convergence is not decoration here. The standard counterexample sets
$f(n,n) = 1$, $f(n,n+1) = -1$, and $f = 0$ elsewhere: summing along rows
gives $0$ every time, hence $0$; summing along columns gives $1$ then $0$
forever, hence $1$. Nothing is infinite or badly behaved --- every row and
every column is a finite sum --- and still the two iterated sums differ,
because the family is not absolutely summable. Rudin's Exercise 3.25 asks for
this phenomenon in a different guise.

Chapter 8 will need this theorem to multiply and rearrange power series, and
Chapter 11 will need its integral analogue.
\end{deeper}

\sitem{S3.13}{Theorem}
\emph{(Means, and the ratio-to-root corollary.) Let $a_n \to a$. Then}
\begin{enumerate}[label=(\alph*), leftmargin=2.2em, itemsep=2pt, topsep=3pt]
  \item $\dfrac{a_1+\cdots+a_n}{n} \to a$ \emph{(Ces\`aro);}
  \item \emph{if moreover $a_n > 0$, then}
        $\sqrt[n]{a_1a_2\cdots a_n} \to a$;
  \item \emph{consequently, if $b_n>0$ and $b_{n+1}/b_n \to a$, then}
        $\sqrt[n]{b_n} \to a$.
\end{enumerate}
\provenance{ECNU \cjk{总练习题}{review} 2 \#3--\#4 (bk.\ p.\,39); USTC
Ch.\ 1 suppl.\ \#6--\#11 (bk.\ pp.\,38--39).}
\begin{proof}
(a) Given $\varepsilon>0$ choose $N$ with $|a_k - a| < \varepsilon/2$ for
$k>N$. For $n>N$,
\[ \Bigl|\frac{1}{n}\sum_{k\le n}a_k - a\Bigr|
   \le \frac{1}{n}\sum_{k \le N}|a_k-a| + \frac{\varepsilon}{2}, \]
and the first term tends to $0$ as $n\to\infty$ with $N$ fixed, so it is
below $\varepsilon/2$ eventually.

(b) If $a>0$, apply (a) to $\log a_n$ and exponentiate, using continuity of
$\exp$. If $a=0$, given $\varepsilon>0$ we have $0<a_k<\varepsilon$ for
$k>N$, and the geometric mean is at most
$\bigl(\prod_{k\le N}a_k\bigr)^{1/n}\varepsilon^{(n-N)/n} \to \varepsilon$.

(c) Put $a_1 = b_1$ and $a_k = b_k/b_{k-1}$ for $k \ge 2$, so that
$a_1a_2\cdots a_n = b_n$ and $a_n \to a$. Now (b) gives
$\sqrt[n]{b_n} \to a$.
\end{proof}

\begin{deeper}[The corollary that explains Rudin's 3.37]
Part (c) is the exact statement behind Rudin's Theorem 3.37,
\[ \liminf \frac{c_{n+1}}{c_n} \le \liminf \sqrt[n]{c_n}
   \le \limsup \sqrt[n]{c_n} \le \limsup \frac{c_{n+1}}{c_n}, \]
whose corollary he states as ``the root test is stronger than the ratio
test.'' Read through (c): the root test sees a \emph{geometric mean} of the
ratios, and geometric means are gentler than the sequence they average ---
they converge whenever it does, and sometimes when it does not. Rudin's
Example 3.35, where the ratios oscillate between $2$ and $1/8$ while the
roots converge, is exactly a case where the mean is tamer than its terms.

Applying (c) to $b_n = n^n/n!$ gives $\lim n/\sqrt[n]{n!} = e$ in one line,
and (a) to $\sqrt[n]{n}$ gives
$\lim \bigl(1+\sqrt2+\sqrt[3]{3}+\cdots+\sqrt[n]{n}\bigr)/n = 1$. Rudin has
the Ces\`aro mean only as Exercise 3.14(a) and never states the geometric
version.
\end{deeper}

\begin{pitfall}[One theorem you will look for and not find]
The Stolz--Ces\`aro theorem --- the discrete l'Hospital rule, that
$\lim (a_n - a_{n-1})/(b_n - b_{n-1}) = L$ forces $\lim a_n/b_n = L$ for
$(b_n)$ strictly increasing and unbounded --- appears \textbf{in none of the
four comparison texts}, and not in Rudin. All four give only the Ces\`aro
special cases of S3.13.

The continuous analogue does appear, as Yu Pin's problem A9$^*$ (p.\,127),
where it is attributed to Cauchy rather than to Stolz; it is exercise S4.3 in
the Chapter 4 appendix. If you want the discrete theorem you will have to
prove it, and S3.13(a) is the case $b_n = n$.
\end{pitfall}
